\begin{table}[ht]
\centering
\renewcommand{\arraystretch}{1.3} % Adjust row height for readability
\begin{tabular}{|p{5cm}|p{5.5cm}|p{5.5cm}|}
\hline
\cellcolor[gray]{0.8}\textbf{Attribute} & \cellcolor[gray]{0.8}\textbf{Option A} & \cellcolor[gray]{0.8}\textbf{Option B} \\ 
\hline
\cellcolor[gray]{0.9}\textbf{What does computer programme do?} 
& Predicts the risk of car theft and highlights areas of particular risk. 
& Identifies individuals who are more likely to commit car theft.   \\ 
\hline
\cellcolor[gray]{0.9}\textbf{Who makes the decision?} 
& The computer program makes a decision that can be overruled by the local police authority.
& The computer program makes a decision that can be overruled by the local police authority. \\ 
\hline
\cellcolor[gray]{0.9}\textbf{Who is performing better on the task?} 
& Computer programs make fewer mistakes than humans.   
& Computer programs and humans make about the same number of mistakes.  \\ 
\hline
\cellcolor[gray]{0.9}\textbf{Who develops the program?} 
& Software engineers of the Ministry of Internal Affairs 
& Software engineers of the Ministry of Internal Affairs  \\ 
\hline
\cellcolor[gray]{0.9}\textbf{Who understands decisions?} 
& Only a developer/software engineer can trace the logic of decision making. 
& The program allows the expert group under the Ministry of Internal Affairs to track the logic of decision-making. \\ 
\hline
\cellcolor[gray]{0.9}\textbf{Who is responsible?} 
& The local police authority that uses this program. 
& Computer program developer/software engineer.     \\ 
\hline
\end{tabular}
\caption{Layout of a Sample Conjoint Task (Originally in Finnish, English Translation)}
\label{tab:conjoint-task-layout3}
\end{table}
\flushleft
\normalsize